\documentclass[11pt]{article}
\newcommand{\ListaTitulo}{Lista 04 --- DCA: Modelo, ANOVA e Número de Réplicas}
% Preâmbulo compartilhado — MATD48, Listas de Exercícios 2026
% Usado por Lista{NN}.tex e Gabarito{NN}.tex via % Preâmbulo compartilhado — MATD48, Listas de Exercícios 2026
% Usado por Lista{NN}.tex e Gabarito{NN}.tex via % Preâmbulo compartilhado — MATD48, Listas de Exercícios 2026
% Usado por Lista{NN}.tex e Gabarito{NN}.tex via \input{preamble}

\usepackage[utf8]{inputenc}
\usepackage[T1]{fontenc}
\usepackage[brazil]{babel}
\usepackage{fullpage}
\usepackage{amsmath,amssymb}
\usepackage{enumitem}
\usepackage{xcolor}
\usepackage{fancyhdr}
\usepackage{titlesec}
\usepackage{listings}
\usepackage{hyperref}
\usepackage{booktabs}
\usepackage{graphicx}

\definecolor{matd48azul}{HTML}{0B4F6C}
\definecolor{matd48verde}{HTML}{2E8B57}

\titleformat{\section}{\normalfont\Large\bfseries\color{matd48azul}}{\thesection}{1em}{}
\titleformat{\subsection}{\normalfont\large\bfseries\color{matd48azul}}{\thesubsection}{1em}{}

\providecommand{\ListaTitulo}{Lista de Exercícios}

\setlength{\headheight}{14pt}
\pagestyle{fancy}
\fancyhf{}
\lhead{\small MATD48 --- Planejamento de Experimentos A}
\rhead{\small \ListaTitulo}
\cfoot{\thepage}
\renewcommand{\headrulewidth}{0.4pt}

\lstset{
  basicstyle=\ttfamily\small,
  breaklines=true,
  frame=single,
  columns=fullflexible,
  backgroundcolor=\color{gray!8},
  commentstyle=\color{matd48verde},
  keywordstyle=\color{matd48azul}\bfseries,
  language=R,
  extendedchars=true,
  literate=%
    {á}{{\'a}}1 {à}{{\`a}}1 {â}{{\^a}}1 {ã}{{\~a}}1
    {é}{{\'e}}1 {ê}{{\^e}}1 {í}{{\'i}}1 {ó}{{\'o}}1
    {ô}{{\^o}}1 {õ}{{\~o}}1 {ú}{{\'u}}1 {ç}{{\c c}}1
    {Á}{{\'A}}1 {À}{{\`A}}1 {Â}{{\^A}}1 {Ã}{{\~A}}1
    {É}{{\'E}}1 {Ê}{{\^E}}1 {Í}{{\'I}}1 {Ó}{{\'O}}1
    {Ô}{{\^O}}1 {Õ}{{\~O}}1 {Ú}{{\'U}}1 {Ç}{{\c C}}1
}

\hypersetup{colorlinks=true, linkcolor=matd48azul, urlcolor=matd48azul, citecolor=matd48azul}

\newcommand{\cabecalho}[2]{%
  \begin{center}
    {\Large\bfseries\color{matd48azul} #1}\\[0.3em]
    {\large #2}\\[0.3em]
    \rule{\linewidth}{0.6pt}
  \end{center}
  \vspace{0.5em}
}

\newenvironment{questao}[1]{\vspace{1em}\noindent\textbf{Questão #1.}\hspace{0.4em}}{\par}
\newenvironment{solucao}{\vspace{0.3em}\noindent\textit{Solução.}\hspace{0.4em}\color{black!85}}{\par\vspace{0.5em}}


\usepackage[utf8]{inputenc}
\usepackage[T1]{fontenc}
\usepackage[brazil]{babel}
\usepackage{fullpage}
\usepackage{amsmath,amssymb}
\usepackage{enumitem}
\usepackage{xcolor}
\usepackage{fancyhdr}
\usepackage{titlesec}
\usepackage{listings}
\usepackage{hyperref}
\usepackage{booktabs}
\usepackage{graphicx}

\definecolor{matd48azul}{HTML}{0B4F6C}
\definecolor{matd48verde}{HTML}{2E8B57}

\titleformat{\section}{\normalfont\Large\bfseries\color{matd48azul}}{\thesection}{1em}{}
\titleformat{\subsection}{\normalfont\large\bfseries\color{matd48azul}}{\thesubsection}{1em}{}

\providecommand{\ListaTitulo}{Lista de Exercícios}

\setlength{\headheight}{14pt}
\pagestyle{fancy}
\fancyhf{}
\lhead{\small MATD48 --- Planejamento de Experimentos A}
\rhead{\small \ListaTitulo}
\cfoot{\thepage}
\renewcommand{\headrulewidth}{0.4pt}

\lstset{
  basicstyle=\ttfamily\small,
  breaklines=true,
  frame=single,
  columns=fullflexible,
  backgroundcolor=\color{gray!8},
  commentstyle=\color{matd48verde},
  keywordstyle=\color{matd48azul}\bfseries,
  language=R,
  extendedchars=true,
  literate=%
    {á}{{\'a}}1 {à}{{\`a}}1 {â}{{\^a}}1 {ã}{{\~a}}1
    {é}{{\'e}}1 {ê}{{\^e}}1 {í}{{\'i}}1 {ó}{{\'o}}1
    {ô}{{\^o}}1 {õ}{{\~o}}1 {ú}{{\'u}}1 {ç}{{\c c}}1
    {Á}{{\'A}}1 {À}{{\`A}}1 {Â}{{\^A}}1 {Ã}{{\~A}}1
    {É}{{\'E}}1 {Ê}{{\^E}}1 {Í}{{\'I}}1 {Ó}{{\'O}}1
    {Ô}{{\^O}}1 {Õ}{{\~O}}1 {Ú}{{\'U}}1 {Ç}{{\c C}}1
}

\hypersetup{colorlinks=true, linkcolor=matd48azul, urlcolor=matd48azul, citecolor=matd48azul}

\newcommand{\cabecalho}[2]{%
  \begin{center}
    {\Large\bfseries\color{matd48azul} #1}\\[0.3em]
    {\large #2}\\[0.3em]
    \rule{\linewidth}{0.6pt}
  \end{center}
  \vspace{0.5em}
}

\newenvironment{questao}[1]{\vspace{1em}\noindent\textbf{Questão #1.}\hspace{0.4em}}{\par}
\newenvironment{solucao}{\vspace{0.3em}\noindent\textit{Solução.}\hspace{0.4em}\color{black!85}}{\par\vspace{0.5em}}


\usepackage[utf8]{inputenc}
\usepackage[T1]{fontenc}
\usepackage[brazil]{babel}
\usepackage{fullpage}
\usepackage{amsmath,amssymb}
\usepackage{enumitem}
\usepackage{xcolor}
\usepackage{fancyhdr}
\usepackage{titlesec}
\usepackage{listings}
\usepackage{hyperref}
\usepackage{booktabs}
\usepackage{graphicx}

\definecolor{matd48azul}{HTML}{0B4F6C}
\definecolor{matd48verde}{HTML}{2E8B57}

\titleformat{\section}{\normalfont\Large\bfseries\color{matd48azul}}{\thesection}{1em}{}
\titleformat{\subsection}{\normalfont\large\bfseries\color{matd48azul}}{\thesubsection}{1em}{}

\providecommand{\ListaTitulo}{Lista de Exercícios}

\setlength{\headheight}{14pt}
\pagestyle{fancy}
\fancyhf{}
\lhead{\small MATD48 --- Planejamento de Experimentos A}
\rhead{\small \ListaTitulo}
\cfoot{\thepage}
\renewcommand{\headrulewidth}{0.4pt}

\lstset{
  basicstyle=\ttfamily\small,
  breaklines=true,
  frame=single,
  columns=fullflexible,
  backgroundcolor=\color{gray!8},
  commentstyle=\color{matd48verde},
  keywordstyle=\color{matd48azul}\bfseries,
  language=R,
  extendedchars=true,
  literate=%
    {á}{{\'a}}1 {à}{{\`a}}1 {â}{{\^a}}1 {ã}{{\~a}}1
    {é}{{\'e}}1 {ê}{{\^e}}1 {í}{{\'i}}1 {ó}{{\'o}}1
    {ô}{{\^o}}1 {õ}{{\~o}}1 {ú}{{\'u}}1 {ç}{{\c c}}1
    {Á}{{\'A}}1 {À}{{\`A}}1 {Â}{{\^A}}1 {Ã}{{\~A}}1
    {É}{{\'E}}1 {Ê}{{\^E}}1 {Í}{{\'I}}1 {Ó}{{\'O}}1
    {Ô}{{\^O}}1 {Õ}{{\~O}}1 {Ú}{{\'U}}1 {Ç}{{\c C}}1
}

\hypersetup{colorlinks=true, linkcolor=matd48azul, urlcolor=matd48azul, citecolor=matd48azul}

\newcommand{\cabecalho}[2]{%
  \begin{center}
    {\Large\bfseries\color{matd48azul} #1}\\[0.3em]
    {\large #2}\\[0.3em]
    \rule{\linewidth}{0.6pt}
  \end{center}
  \vspace{0.5em}
}

\newenvironment{questao}[1]{\vspace{1em}\noindent\textbf{Questão #1.}\hspace{0.4em}}{\par}
\newenvironment{solucao}{\vspace{0.3em}\noindent\textit{Solução.}\hspace{0.4em}\color{black!85}}{\par\vspace{0.5em}}


\begin{document}

\cabecalho{Lista de Exercícios 04}{Delineamento completamente aleatorizado: modelo, ANOVA, intervalos de confiança e tamanho amostral (Capítulo 3 / Aula 04)}

\noindent \textit{Entrega: até a próxima aula. Justifique todas as respostas --- nestas listas, a
justificativa vale mais que a resposta final. O gabarito completo está em \texttt{Gabarito04.pdf}.}

\begin{questao}{1} \textbf{(Psicologia --- ANOVA a partir de somas de quadrados)}

Uma psicóloga compara três técnicas de relaxamento (respiração guiada, meditação breve e
alongamento) quanto à redução do nível de cortisol salivar (em nmol/L) após uma tarefa
estressante. Ela sorteia 30 voluntários, 10 por técnica, e obtém $SQ_{\text{Trat}} = 148{,}2$ e
$SQ_{\text{Total}} = 402{,}6$.

\begin{enumerate}[label=(\alph*)]
  \item Complete a tabela de ANOVA (fonte, g.l., SQ, QM, $F$).
  \item Ao nível $\alpha = 0{,}05$, com $F_{2,27;0{,}05} \approx 3{,}35$, as três técnicas diferem
    quanto à redução média de cortisol? Justifique com base no valor de $F$ calculado.
  \item O que exatamente essa conclusão permite afirmar --- e o que ela \emph{não} permite
    afirmar --- sobre quais pares de técnicas diferem entre si?
\end{enumerate}
\end{questao}

\begin{questao}{2} \textbf{(Agricultura --- intervalos de confiança)}

Um experimento com quatro doses de um fertilizante nitrogenado (0, 40, 80 e 120 kg/ha) sobre a
produtividade do milho (sc/ha) usa 8 parcelas por dose ($N=32$). O ajuste da ANOVA fornece
$MS_E = 18{,}4$ com $28$ graus de liberdade, e as médias observadas são $\bar Y_{0} = 62{,}1$,
$\bar Y_{40} = 68{,}7$, $\bar Y_{80} = 74{,}3$ e $\bar Y_{120} = 73{,}5$. Use
$t_{28;0{,}025} \approx 2{,}048$.

\begin{enumerate}[label=(\alph*)]
  \item Construa um IC de 95\% para a produtividade média da dose de 80 kg/ha.
  \item Construa um IC de 95\% para a diferença de produtividade entre as doses de 80 e 40 kg/ha.
  \item Por que os dois intervalos acima usam o mesmo $MS_E = 18{,}4$, em vez de cada dose estimar
    sua própria variância a partir de suas 8 parcelas?
\end{enumerate}
\end{questao}

\begin{questao}{3} \textbf{(Ciência de dados --- notação matricial do DCA)}

Considere um DCA com $t = 3$ tratamentos e $n_1 = 4$, $n_2 = 3$, $n_3 = 5$ unidades ($N=12$),
usando a parametrização de médias de célula $\mathbf{Y} = \mathbf{X}\boldsymbol\beta +
\boldsymbol\varepsilon$, com $\mathbf{X}$ a matriz $N\times t$ de indicadoras de tratamento e
$\boldsymbol\beta = (\mu_1,\mu_2,\mu_3)'$.

\begin{enumerate}[label=(\alph*)]
  \item Dê as dimensões e o conteúdo explícito de $\mathbf{X}'\mathbf{X}$ e de
    $(\mathbf{X}'\mathbf{X})^{-1}$.
  \item Explique por que, nesta parametrização (diferente da parametrização $\mu + \tau_i$ com
    restrição de soma zero), \emph{todo} $\mu_i$ é trivialmente estimável.
  \item Usando $\mathrm{Cov}(\hat{\boldsymbol\beta}) = \sigma^2(\mathbf{X}'\mathbf{X})^{-1}$, dê a
    expressão de $\mathrm{Var}(\hat\mu_2 - \hat\mu_3)$ em função de $\sigma^2$.
\end{enumerate}
\end{questao}

\begin{questao}{4} \textbf{(Dedução --- $\hat{\boldsymbol\beta}$ é BLUE no DCA)}

Seja $\mathbf{X}$ a matriz $N\times t$ de indicadoras de tratamento (uma coluna por tratamento,
sem intercepto) do DCA, de posto coluna completo $t$, e $\mathbf{Y} = \mathbf{X}\boldsymbol\beta +
\boldsymbol\varepsilon$ com $E[\boldsymbol\varepsilon] = \mathbf{0}$ e
$\mathrm{Cov}(\boldsymbol\varepsilon) = \sigma^2\mathbf{I}_N$ (sem exigir normalidade).

\begin{enumerate}[label=(\alph*)]
  \item Mostre que $\hat{\boldsymbol\beta} = (\mathbf{X}'\mathbf{X})^{-1}\mathbf{X}'\mathbf{Y}$ é
    não viesado para $\boldsymbol\beta$.
  \item Mostre que $\mathrm{Cov}(\hat{\boldsymbol\beta}) = \sigma^2(\mathbf{X}'\mathbf{X})^{-1}$ e
    que essa matriz é diagonal, com $i$-ésimo elemento $\sigma^2/n_i$.
  \item Enuncie o Teorema de Gauss-Markov e explique por que suas hipóteses estão satisfeitas
    aqui, concluindo que $\hat\mu_i = \bar Y_{i\cdot}$ é o BLUE de $\mu_i$ para cada $i$.
\end{enumerate}
\end{questao}

\begin{questao}{5} \textbf{(Número de réplicas --- ciência de dados)}

Uma equipe de engenharia testa três estratégias de \emph{cache} (nenhuma, LRU, LFU) quanto ao
tempo de carregamento de uma página (ms). Um piloto pequeno sugere médias aproximadas de
$320$, $290$ e $270$ ms e um desvio-padrão comum $\sigma \approx 55$ ms. A tabela abaixo mostra o
poder de um teste ANOVA a $\alpha = 0{,}05$, para $k=3$ grupos e $f = 0{,}30$ (o tamanho de efeito
de Cohen implícito nessas médias piloto), em função do número de réplicas $n$ por grupo:

\begin{center}
\begin{tabular}{lcccccc}
\toprule
$n$ & 15 & 20 & 30 & 40 & 50 & 60 \\
\midrule
Poder & 0{,}38 & 0{,}50 & 0{,}69 & 0{,}80 & 0{,}88 & 0{,}92 \\
\bottomrule
\end{tabular}
\end{center}

\begin{enumerate}[label=(\alph*)]
  \item Qual o menor $n$ por grupo, entre os valores tabelados, que garante poder de pelo menos
    $0{,}80$? Quantas observações totais isso representa?
  \item Se o orçamento só permitir $n=20$ por grupo, o que isso implica sobre a chance de a
    equipe \emph{não} detectar uma diferença real do tamanho assumido no piloto?
  \item O valor de $\sigma$ usado veio de um piloto pequeno e é, portanto, incerto. Que direção de
    erro nessa estimativa (superestimar ou subestimar $\sigma$) tornaria o planejamento de $n$
    baseado nesta tabela otimista demais (poder real menor que o nominal)?
\end{enumerate}
\end{questao}

\end{document}
